\documentclass[11pt,a4paper]{article}
\usepackage[utf8]{inputenc}
\usepackage[ngerman]{babel}
\usepackage[T1]{fontenc}
\usepackage{xcolor}
\usepackage{hyperref}
\usepackage{tabularx}

\title{Projektplan \& Status: Echtzeitsysteme (SS26)}
\author{Danielou Mounsande}
\date{Stand: 05. Mai 2026}

\begin{document}

\maketitle

\section{Aktueller Status (KW 19 - 05. Mai)}
Das Projekt befindet sich exakt im Zeitplan. Der \textbf{Meilenstein Minimalprojekt} (2 Tasks, Round-Robin) ist de facto erreicht.

\subsection{Abgeschlossen (Basis-Funktionen)}
\begin{itemize}
    \item \textbf{Setup:} Projekt \& System aufgesetzt.
    \item \textbf{Task-Implementierung:} TCB-Struktur definiert und Task-Erstellung (Dummy Stack Frame) implementiert.
    \item \textbf{Context Switch:} Hybrider Wechsel (PendSV Assembler + C-Logik). Sichern von R4-R11 und Hardware-Stacking funktionieren fehlerfrei.
    \item \textbf{Scheduler RR:} Einfaches präemptives Round-Robin Scheduling via SysTick (Zeitscheiben) ist aktiv.
\end{itemize}

\section{Der Weg nach vorn (Ab KW 20)}
Gemäß dem Zeitplan von Prof. Schmidt aus \texttt{02\_Tasks\_SS26.pdf} beginnt nun die Phase der \textbf{Kernel Objekte}.

\subsection{Nächster Schritt (Priorität 1): Blocking-State \& Delay}
Aktuell nutzen Tasks \textit{Busy-Waiting} (for-Schleifen). Das Skript (Folie 26) fordert zwingend einen \textbf{Blocking-State}.
\begin{itemize}
    \item \textbf{To-Do:} TCB um \texttt{u32TicksToWait} erweitern.
    \item \textbf{To-Do:} \texttt{Kernel\_Delay(ticks)} implementieren (Task in Blocked-State versetzen, Yield anfordern).
    \item \textbf{To-Do:} SysTick-Handler anpassen, um wartende Tasks zu aktualisieren.
\end{itemize}

\subsection{Kommende Kernel Objekte (Zeitplan)}
\begin{table}[h]
    \centering
    \begin{tabularx}{\textwidth}{|l|l|X|}
    \hline
    \textbf{Woche} & \textbf{Thema} & \textbf{Ziel} \\ \hline
    KW 20 & Scheduler Update & Verfeinerung Präemptives RR \& Task States. \\ \hline
    KW 21 & Synchronisation & Semaphoren \& Mutexe implementieren. \\ \hline
    KW 22 & Kommunikation & Message Queues einbauen. \\ \hline
    KW 23 & \textbf{Meilenstein} & \textbf{Kernel-Objekte abgeschlossen.} \\ \hline
    \end{tabularx}
\end{table}

\subsection{Verifizierung \& Dokumentation (KW 24+)}
\begin{itemize}
    \item \textbf{Instrumentierung:} Integration von Segger SystemView zur visuellen Analyse des Schedulers.
    \item \textbf{Test-Pipeline:} Setup der TeSSLa Pipeline und Spezifizierung der Tests.
    \item \textbf{Abschluss:} Dokumentation verfassen und Puffer-Zeit nutzen.
\end{itemize}

\end{document}
